% !TeX program = xelatex
\documentclass[11pt,a4paper]{article}

\usepackage{fontspec}
\usepackage{xeCJK}
\usepackage{amsmath,amssymb}
\usepackage{array}
\usepackage[table]{xcolor}
\usepackage{enumitem}
\usepackage[
  top=2.15cm,
  bottom=2.15cm,
  left=2.2cm,
  right=2.2cm
]{geometry}

\setmainfont{Arial}
\setsansfont{Arial}
\setCJKmainfont{Microsoft YaHei}
\setCJKsansfont{Microsoft YaHei}
\newCJKfontfamily\cnhei{SimHei}

\pagestyle{empty}
\setlength{\parindent}{0pt}
\setlength{\parskip}{0.35em}
\linespread{1.16}
\setlist[enumerate]{label=\arabic*.,leftmargin=2em,itemsep=0.2em,topsep=0.2em}
\setlength{\abovedisplayskip}{0.55em}
\setlength{\belowdisplayskip}{0.55em}

\newcommand{\examtitle}[1]{%
  {\cnhei\fontsize{20}{24}\selectfont #1}\par
  \vspace{0.35em}
  {\color{gray!25}\hrule height 0.5pt}
  \vspace{0.6em}
}

\newcommand{\studentline}{%
  姓名：\underline{\makebox[3.0cm]{}}\quad
  学号：\underline{\makebox[3.2cm]{}}\par
  \vspace{0.7em}
  {\color{gray!25}\hrule height 0.5pt}
  \vspace{1.0em}
}

\newcommand{\gradetable}{%
  {\cnhei 评分表(每题10分, 共100分)}\par\vspace{0.45em}
  {\renewcommand{\arraystretch}{1.75}
  \setlength{\tabcolsep}{4pt}
  \arrayrulecolor{gray!45}
  \begin{tabular*}{\textwidth}{@{\extracolsep{\fill}}|c|*{10}{c|}c|}
    \hline
    \rule{0pt}{2.0em}{\cnhei 题号} & {\cnhei 1} & {\cnhei 2} & {\cnhei 3} & {\cnhei 4} & {\cnhei 5} & {\cnhei 6} & {\cnhei 7} & {\cnhei 8} & {\cnhei 9} & {\cnhei 10} & {\cnhei 总分} \\
    \hline
    \rule{0pt}{2.0em}{\cnhei 得分} & & & & & & & & & & & \\
    \hline
  \end{tabular*}
  \arrayrulecolor{black}}\par
  \vspace{1.0em}
  {\color{gray!25}\hrule height 0.5pt}
  \vspace{0.8em}
}

\newcommand{\sectiontitle}[1]{%
  \vspace{0.9em}
  {\cnhei #1}\par
  \vspace{0.55em}
}

\newcommand{\question}[1]{\par\noindent{\cnhei #1.}\quad}
\newcommand{\answerspace}[1]{\par\vspace*{#1}}

\begin{document}

\examtitle{《新工科数学分析》期中测试卷II}

{\small 考试时间：80 分钟\quad 满分：100 分}\par
\vspace{0.7em}

\studentline
\gradetable


\question{1}设 \(z=e^{xy}+x^2y^3,\ x=r\cos\theta,\ y=r\sin\theta\)。求
\(\dfrac{\partial z}{\partial r},\ \dfrac{\partial z}{\partial \theta}\)。
\answerspace{4.2cm}




\question{2}设方程 \(x^2+y^2+z^2+xy-xz+2yz-4=0\) 确定隐函数
\(z=z(x,y)\)，求
\(\dfrac{\partial z}{\partial x},\ \dfrac{\partial z}{\partial y}\)。
\answerspace{4.1cm}

\question{3}设方程组
\[
\begin{cases}
xu+yv-u^2=1,\\
x^2+y^2+u+v=3
\end{cases}
\]
求
\(\dfrac{\partial u}{\partial y},\ \dfrac{\partial v}{\partial y}\)。
\answerspace{2.3cm}

\newpage

\question{4}求函数 \(f(x,y)=x^3+y^3-3xy\) 的极值。
\answerspace{4cm}

\question{5}求抛物面 \(z=x^2+y^2\) 上到点 \((1,2,0)\) 距离最近的点。
\answerspace{4cm}


\question{6}用拉格朗日乘子法求函数
\(f(x,y,z)=x^2+y^2+z^2\) 在条件 \(x+y+z=6\) 下的最小值。
\answerspace{4cm}


\question{7}计算二重积分 \(\displaystyle \iint_D xy\,dA\)，其中
\(D=\{(x,y)\mid x\ge0,\ y\ge0,\ x+y\le2\}\)。
\answerspace{4cm}


\question{8}计算三重积分 \(\displaystyle \iiint_\Omega (x+y+z)\,dV\)，其中
\(\Omega=\{(x,y,z)\mid 0\le x\le1,\ 0\le y\le2,\ 0\le z\le3\}\)。
\answerspace{4cm}

\newpage

\question{9}设二维稳态温度场 \(u(x,y)=e^x\cos y\)。
\begin{enumerate}
\item 计算 \(u_x,\ u_y,\ u_{xx},\ u_{yy}\)。
\item 验证 \(u\) 是否满足 Laplace 方程 \(u_{xx}+u_{yy}=0\)。
\item 计算点 \((0,\pi/3)\) 处的梯度 \(\nabla u\)。
\item 求该点处温度增长最快的方向及最大变化率。
\end{enumerate}
\answerspace{6.6cm}

\question{10}设函数 \(u(x,y,t)=e^{-2t}\sin x\sin y\)。
\begin{enumerate}
\item 求 \(u_t,\ u_{xx},\ u_{yy}\)。
\item 验证 \(u\) 是否满足二维热传导方程 \(u_t=u_{xx}+u_{yy}\)。
\item 求时刻 \(t=0\) 时温度场在点
\(\left(\dfrac{\pi}{4},\dfrac{\pi}{3}\right)\) 处沿方向
\(\mathbf l=\left(\dfrac35,\dfrac45\right)\) 的方向导数。
\end{enumerate}
\answerspace{6.6cm}

\end{document}
